% 由 scripts/publication/tables.py 自动生成，请勿手改。
% 需要宏包：booktabs（\usepackage{booktabs}）
\begin{table}[htbp]
  \centering
  \caption{Step 12 主结果：A–D 方案（256/512 通道）}
  \label{tab:step12-main}
  \begin{tabular}{lrrrrrrrl}
    \toprule
    方案 & 平均损耗 (dB) & P95 (dB) & 最大损耗 (dB) & 未布通 & 唯一交叉 & 间距违规 & 运行 (s) & 完整连接 \\
    \midrule
    A & 5.2786 & 6.1318 & 6.3333 & 0 & 9\,468 & 393 & 0.7 & 是 \\
    B & 4.3360 & 5.1613 & 5.3638 & 0 & 9\,468 & 33 & 12.9 & 是 \\
    C & 5.2772 & 6.1310 & 6.3222 & 0 & 9\,484 & 380 & 23.2 & 是 \\
    D & 4.3394 & 5.1587 & 5.3462 & 0 & 9\,568 & 39 & 38.6 & 是 \\
    \midrule
    A & 5.5162 & 6.3510 & 6.5773 & 0 & 43\,393 & 692 & 1.8 & 是 \\
    B & 4.5582 & 5.3886 & 5.6217 & 2 & 43\,306 & 522 & 40.9 & 否 \\
    C & 5.5146 & 6.3484 & 6.5727 & 0 & 43\,320 & 703 & 94.0 & 是 \\
    D & 4.5593 & 5.3886 & 5.6129 & 0 & 43\,170 & 516 & 159.2 & 是 \\
    \bottomrule
  \end{tabular}
  \par\vspace{2pt}
  \begin{flushleft}\footnotesize
  A=原版 R5；B=统一 R6（参数对照，非算法贡献）；C=候选轨道优化（R5）；D=自适应半径（R5/R6）。 \\
  512 的 B 有 2 条未布通，其损耗均值与完整连接的方案不可直接比较；排名只在完整连接的方案间进行。 \\
  最优值加粗仅在同一规模、同一完整性条件的行内比较。 \\
  损耗为解析物理统计口径（交叉事件按每根波导各计一次）；与原版 calc\_index 统计口径不同，两者不混用。 \\
  \end{flushleft}
\end{table}
